\documentclass{article}
\usepackage{graphicx} % Required for inserting images

\title{Investment Appraisal}
\author{Dalya Al Sallo, Emily Latchem, Samhal Guesh, Tahera Hussain}
\date{April 2026}

\begin{document}

\maketitle

\section{Investment Appraisal}

Investment appraisal is the process used by businesses and investors to evaluate whether a project is financially worthwhile. It involves comparing initial costs with expected future cash flows to determine profitability and risk. This forms a key part of capital budgeting decisions (Brealey, Myers \& Allen, 2020).

A fundamental concept is the time value of money, which states that money today is worth more than in the future due to its earning potential. Therefore, future cash flows are discounted to present value to allow fair comparison (Ross, Westerfield \& Jaffe, 2019).

Common techniques include Net Present Value (NPV) and Internal Rate of Return (IRR). Other methods such as Payback Period and Profitability Index assess liquidity and efficiency (Umeorah, 2024).

\section{UK Government Support}

Government schemes can improve investment outcomes. The Annual Investment Allowance allows businesses to deduct up to £1 million of qualifying capital expenditure from profits, improving NPV (HM Revenue \& Customs, 2023a).

R\&D tax relief reduces the cost of innovation and can increase returns on investment (HM Revenue \& Customs, 2023b). Enterprise Investment Schemes provide tax incentives to investors, encouraging investment and lowering the cost of capital (HM Revenue \& Customs, 2023c). Grants, particularly for energy efficiency and net-zero projects, reduce initial costs and improve project feasibility.

\section{Library Overview}

The library calculates key investment appraisal metrics from cash flow data and provides simple interpretation of results.

\begin{verbatim}
>>> import investment_appraisal as ia
>>> ia.npv(0.12, [-100000, 70000, 85000, 65000])
7527.20
\end{verbatim}

A positive NPV indicates the investment adds value.

\begin{verbatim}
>>> ia.irr([-100000, 70000, 85000, 65000])
0.53
\end{verbatim}

This suggests a return of 53\%, well above the required rate.

\subsection*{Available Methods}
Present Value, Future Value, NPV, IRR, Payback Period, Discounted Payback, Profitability Index, ARR, Annuity PV, Perpetuity, Growing Perpetuity.

\section{Statement of Need}

This library combines multiple investment appraisal techniques into a single, accessible tool. While similar functionality exists in libraries such as NumPy Financial, this library focuses on simplicity and interpretation.

It allows users to evaluate projects, compare alternatives, and understand key financial concepts such as time value of money.

\section{Conclusion}

Investment appraisal is essential for evaluating project viability and supporting capital budgeting decisions. This library provides a simple and practical tool for applying key methods such as NPV and IRR.

Future improvements could include tools for comparing multiple investment projects simultaneously. Overall, the library provides a useful foundation for both financial analysis and learning.

\section*{References}

Brealey, R. A., Myers, S. C. \& Allen, F. (2020) \textit{Principles of Corporate Finance}.

Ross, S. A., Westerfield, R. W. \& Jaffe, J. (2019) \textit{Corporate Finance}.

Umeorah, N. (2024) Lecture Notes. Cardiff University. Note - not available to public.

HM Revenue \& Customs (2023a) \textit{Evaluation of the Annual Investment Allowance}. Available at: https://www.gov.uk/

HM Revenue \& Customs (2023b) \textit{Research and Development Tax Relief Guidance}. Available at: https://www.gov.uk/

HM Revenue \& Customs (2023c) \textit{Enterprise Investment Scheme Guidance}. Available at: https://www.gov.uk/

\end{document}